With the deep evolution of enterprise digital transformation and hybrid cloud architectures, the network attack surface of IT infrastructure is expanding exponentially. Traditional static defense systems and manual penetration testing are increasingly inadequate against Advanced Persistent Threats (APTs) characterized by long lifecycles and multi-stage behaviors. In this context, Breach and Attack Simulation (BAS) has emerged as a core pillar for Continuous Threat Exposure Management (CTEM). However, existing open-source automated penetration testing frameworks generally suffer from severe engineering limitations, notably the tight coupling of control flow, network I/O, and penetration business logic in their underlying orchestration engines. When simulating complex, long-chain kill chains, such concurrent execution architectures lacking strict state management are highly prone to memory state pollution, logic deadlocks, and a lack of fault tolerance.

To address these engineering pain points, this thesis designs and implements a lightweight automated penetration testing orchestration engine based on Microkernel Architecture and Finite State Machine (FSM). First, the system introduces the Dependency Inversion Principle (DIP) to construct a physically isolated I/O runtime layer, strictly decoupling complex environmental interactions from core scheduling operations. Second, a data-driven orchestration model based on an externalized state transition matrix is proposed. This degrades penetration tactical operators to stateless pure functions, completely eliminating control fragmentation. Finally, an absolute single-threaded step-by-step orchestration combined with a transactional snapshot persistence mechanism is adopted, achieving highly reliable, lock-free breakpoint continuation.

Based on this orchestration engine, the thesis fully implements eight core attack abilities mapped to the MITRE ATT\&CK framework, encompassing low-level adversarial techniques such as DNS SRV-based domain reconnaissance, COM integrity escape for privilege escalation, kernel driver blinding and ETW blocking, offline credential cracking, and fileless lateral movement. Engineering validation demonstrates that the engine successfully overcomes the state loss dilemma in long-chain automated penetration testing, exhibiting exceptional practical stealth and industrial-grade disaster recovery capabilities. This provides a standardized and highly scalable engineering practice scheme for the control plane architecture evolution of next-generation BAS systems.
